% Part: many-valued-logic
% Chapter: infinite-valued-logics
% Section: introduction

\documentclass[../../../include/open-logic-section]{subfiles}

\begin{document}

\olfileid{mvl}{inf}{int}

\olsection{جان پہچان}

میٹرکس وچ صداقتی قدراں دی گنتی محدود ہونی ضروری نہیں۔
لامتناہی صداقتی قدراں لئی اک صاف چُون $0$ تے~$1$ دے وچکارلے
ناطق عدداں دا سیٹ $V_\infty = [0,1] \cap \Rat$ اے، یعنی
\begin{align*}
    V_\infty & = \Setabs{\frac{n}{m}}{n,m \in \Nat \text{ تے } n\le m}.
\intertext{صداقتی قدراں دے ایس لامتناہی سیٹ نوں ویکھدیاں ایہ ذیلی سیٹ وی اکثر کم آؤندے نیں۔}
V_m & = \Setabs{\frac{n}{m-1}}{n \in \Nat \text{ تے } n\le m}
\intertext{مثلاً $V_5$ اوہ سیٹ اے جیہدیاں $5$ صداقتی قدراں برابر فاصلے اُتے نیں۔}
V_5 & = \{0, \frac{1}{4}, \frac{1}{2}, \frac{3}{4}, 1\}.
\end{align*}
ایہناں صداقتی سیٹاں اُتے بنیاں منطقاں وچ عام طور اُتے صرف
$1$ مقرر ہوندا اے، یعنی $V^+ = \{1\}$۔ دوجے لفظاں وچ
$1$ نوں (پورا) سچ تے $0$ نوں پورا جھوٹ سمجھدے آں، پر
!!{formula}s~$V$ وچ کوئی وی درمیانی قدر لے سکدے نیں۔

فیر وی اسیں $V_{[0,1]} = [0,1]$، یعنی $0$ تے~$1$
دے وچکارلے سارے \emph{حقیقی} عدداں دا سیٹ، یا $[0,1]$ دے
ہور لامتناہی ذیلی سیٹ وی ویکھ سکدے آں۔ ایس صداقتی سیٹ والیاں
منطقاں نوں اکثر \emph{دھندلیاں منطقاں} آکھدے آں۔

\end{document}
